% ============================================================
% Problema 10
% ============================================================
\section{Problema 10: Capa 4: El servidor no responde}

\begin{tcolorbox}[enunciado]
Un usuario puede hacer ping al servidor correctamente. Sin embargo, cuando intenta acceder a una aplicación web mediante \texttt{http://192.168.50.10}, la conexión falla. La PC tiene una dirección IP correcta y el servidor responde a ping. El servidor tiene habilitado el servicio web y la aplicación utiliza el puerto TCP 80.

\begin{enumerate}[label=\arabic*., leftmargin=*]
    \item ¿Qué capa del modelo OSI está involucrada principalmente?
    \item ¿Qué protocolo de transporte se está utilizando?
    \item ¿Qué diferencia existe entre comprobar conectividad con ping y comprobar un puerto TCP?
    \item ¿Qué podría estar bloqueando la conexión?
    \item ¿Cómo comprobarías si el puerto 80 está accesible?
\end{enumerate}
\end{tcolorbox}

\subsection*{Solución}
\begin{enumerate}[label=\arabic*., leftmargin=*]
    \item ¿Qué capa del modelo OSI está involucrada principalmente?
    
    La Capa 4 (Transporte): Porque el enrutamiento básico funciona, pero el problema 
    ocurre al intentar conectar con un puerto TCP específico para levantar un servicio.

    \item ¿Qué protocolo de transporte se está utilizando?
    
    El protocolo TCP, ya que es el encargado de mantener la conexión de los servicios web en el puerto 80.
    
    \item ¿Qué diferencia existe entre comprobar conectividad con ping y comprobar un puerto TCP?
    
    El ping (Capa 3) comprueba si un dispositivo está encendido y responde en la red, y comprobar un puerto TCP (Capa 4) 
    verifica si un servicio está encendido y aceptando conexiones.
    
    \item ¿Qué podría estar bloqueando la conexión?
    
    Un Firewall que esté bloqueando específicamente el tráfico hacia el puerto 80.
    
    \item ¿Cómo comprobarías si el puerto 80 está accesible?
    
    Usando el comando telnet 192.168.50.10 80.
\end{enumerate}

